\newcommand{\CaseID}{camp\_fire}
\newcommand{\CaseTitle}{2018 Camp Fire Landscape Reconstruction}
\newcommand{\EventStart}{8 November 2018 at 14:00 PST}
\newcommand{\SimulationDuration}{36 hours}
\newcommand{\ComparisonTime}{10 November 2018 at 02:00 PST (10:00 UTC)}
\newcommand{\GridDescription}{1457 by 1476 cells at 30 m in EPSG:26910}
\newcommand{\WeatherDescription}{37-band weather rasters; bands 1--25 are selected at one-hour intervals}
\newcommand{\IgnitionDescription}{random ignition sampling from the archived ignition mask}
\newcommand{\ArchiveName}{Camp\_VAL.rar}
\newcommand{\ArchiveSha}{\seqsplit{0b0e8f7ce3bbea4665269ff7f5b88dd373f3d9210b3fa4ddc7916b610f724033}}
\newcommand{\InputCaveat}{No additional active-range anomaly was identified by the structural and statistical checks; upstream plausibility remains unverified because product metadata are absent.}
\newcommand{\CompatibilityNote}{Obsolete domain-description and diagnostic settings were omitted. The earlier spotting choices were expressed as fire-power-proportional generation, empirical transport distance, Eulerian accumulation, and physical ignition.}
